\begin{lemma}
\label{editorial-source-lemma-universally-bijective}
Let $\varphi : R \to S$ be a ring map such that
\begin{enumerate}
\item the kernel of $\varphi$ is locally nilpotent, and
\item $S$ is generated as an $R$-algebra by elements $x$
such that there exist $n > 0$ and a polynomial $P(T) \in R[T]$
whose image in $S[T]$ is $(T - x)^n$.
\end{enumerate}
Then $\Spec(S) \to \Spec(R)$ is a homeomorphism and $R \to S$
induces purely inseparable extensions of residue fields.
Moreover, conditions (1) and (2) remain true on arbitrary base change.
\end{lemma}

\begin{proof}
If $S=0$, the nil-kernel hypothesis
makes $1\in R$ nilpotent, so $R=0$.
Every unital base change of the zero
ring is again zero, and all assertions
follow. Hence assume $S\neq0$.
Retain the original map $\varphi:R\to S$
and its kernel $I$. For every original
generator $x$ retain the given positive
integer $n$ and the entire original
polynomial $P(T)=\sum_{j=0}^d c_jT^j$,
whose image is $(T-x)^n$. Coefficient
comparison gives $d\geq n$,
$c_n-1\in I$, and $c_j\in I$ for
$j>n$. The polynomial
$Q(T)=T^n+\sum_{j=0}^{n-1}c_jT^j$
is a monic annihilating polynomial
for $x$. Indeed the full difference
$$
P(T)-Q(T)=(c_n-1)T^n+
\sum_{j=n+1}^d c_jT^j
$$
has all coefficients in the original
kernel, so $Q$ has the same image
$(T-x)^n$. This retains every
coefficient of $P$ and proves
integrality without presuming that
the given lift itself is monic.
All generators, hence $S$, are
integral over $R$. The integral
injection $R/I\to S$ is surjective
on spectra. Since $I$ is nil,
$V(I)=\Spec(R)$, so the original
spectrum map is surjective; it is
also closed by integrality.

We first prove the global scalar
identity used in the fibre argument.
Write
$$
e_j=\binom nj x^j
=\varphi((-1)^j c_{n-j})\in\varphi(R),
\qquad 1\leq j\leq n.
$$
All signs here are the actual signs
in the given $(T-x)^n$. For
$1\leq h\leq n$, the full Newton
recurrence is
$$
n x^h=
\sum_{j=1}^{h-1}(-1)^{j-1}e_j\,n x^{h-j}
+(-1)^{h-1}h e_h.
$$
It follows directly on substituting
the original $e_j$: Pascal's identity
telescopes to
$\sum_{j=0}^{h-1}(-1)^j\binom nj
=(-1)^{h-1}\binom{n-1}{h-1}$,
and $n\binom{n-1}{h-1}=h\binom nh$.
These two integer identities give
exactly the coefficient $n$ on the
right, with every displayed summand
retained. Induction on $h$, starting
with $n x=e_1$, proves $n x^h\in\varphi(R)$.
Explicit preimages are supplied by
the same recurrence in $R$, taking
$E_j=(-1)^j c_{n-j}$ and
$$
U_h=\sum_{j=1}^{h-1}(-1)^{j-1}E_j U_{h-j}
+(-1)^{h-1}h E_h.
$$
The empty sum for $h=1$ gives
$U_1=E_1$, and $\varphi(U_h)=n x^h$.

Fix a prime integer $p$ and write
the original $n=p^e m$, with
$p\nmid m$, retaining $q=p^e$.
One has the exact integer identity
$$
\gcd\left(n,\binom nq\right)=m.
$$
For a prime $\ell\neq p$ dividing
$n$, the identity
$\binom nq=(n/q)\binom{n-1}{q-1}$
shows its $\ell$-adic valuation is
at least that of $n$, since
$q$ is prime to $\ell$.
For $p$, the factor $n/q=m$
is prime to $p$, and
$\binom{n-1}{q-1}$ is nonzero
modulo $p$. In fact
$(1+T)^{qm}=(1+T^q)^m$ modulo
$p$, so dividing by $1+T$ shows
that its degree-$q-1$ coefficient
is $(-1)^{q-1}=1$ in
$\mathbf F_p$. This also treats
$q=1$. These valuations prove
the displayed greatest common
divisor, including every prime
factor of the original $n$.
Choose integers $A,B$ with
$An+B\binom nq=m$. Then
$$
\varphi(AU_q+BE_q)
=A n x^q+B\binom nq x^q=m x^q.
$$
Thus the source assertion
$m x^{p^e}\in\varphi(R)$ is
indeed true globally, not merely
after taking a residue field.
The recurrence and this integer
identity supply its missing proof.

Now let $\mathfrak p\subset R$
and retain $K=\kappa(\mathfrak p)$
and the original fibre
$C=K\otimes_R S$. It is nonzero
by surjectivity on spectra. Thus
the map from the field $K$ to $C$
is injective. If $K$ has
characteristic zero, each generator
$y=1\otimes x$ satisfies
$n y=-\overline c_{n-1}$.
The original positive $n$ is
invertible in $K$, so
$y=-n^{-1}\overline c_{n-1}$
belongs to the image of $K$.
All generators do, and hence
$K\to C$ is an isomorphism.

If $K$ has characteristic $p>0$,
use the original decomposition
$n=q m$, $q=p^e$, for each
generator separately. The full
fibre polynomial is
$$
(T-y)^n=(T^q-y^q)^m
=\sum_{j=0}^m\binom mj(-1)^j
T^{q(m-j)}y^{qj}.
$$
Its coefficient of $T^{n-q}$
is $-m y^q=\overline c_{n-q}$.
Since $m$ is a unit of $K$,
$y^q=-m^{-1}\overline c_{n-q}$
is in the original image of $K$.
This is the direct fibre formula;
the global preimage proved above
gives the same conclusion before
reducing coefficients.
If $e=0$, use the exponent $1$
in Lemma \ref{editorial-source-lemma-p-ring-map};
otherwise use $e$. In either
case a positive $p$-power of $y$
is in $K$, and the corresponding
integer multiple is zero in
characteristic $p$. Thus that
lemma applies to $K\to C$, whose
kernel is zero. The fibre has
one prime and its residue field
is purely inseparable over $K$.
The same conclusion holds in
characteristic zero by the
already proved isomorphism.
Remark \ref{algebra-remark-fundamental-diagram}
now gives bijectivity on spectra
and the asserted original residue
extensions. A continuous closed
bijection is a homeomorphism.

Under any original $f:R\to R'$,
the image polynomial
$P_f(T)=\sum_j f(c_j)T^j$ has
image $(T-1\otimes x)^n$ in
$(R'\otimes_R S)[T]$, with all
coefficients and the same $n$
retained. These elements generate
the base-changed algebra, so
condition (2) persists. For
condition (1), the original
surjectivity on spectra persists
under base change by the explicit
nonzero residue-field tensor
construction in Lemma
\ref{editorial-source-lemma-2-3-ring-map}. Hence
the new kernel lies in every
prime and is locally nilpotent.
Both hypotheses, and therefore
all conclusions, hold for the
actual base-changed map.

The scalar identities give more
global information for a finite
original generating list
$x_1,\ldots,x_r$ with integers
$n_1,\ldots,n_r$. Put
$N=\prod_{i=1}^r n_i$, taking
$N=1$ for the empty list.
Since $n_i x_i$ has the
specified preimage $U_{i,1}$,
the map $R[1/N]\to S[1/N]$
is surjective, sending the
original fraction
$U_{i,1}/n_i=((N/n_i)U_{i,1})/N$
to $x_i$. Its kernel is exactly
$I[1/N]$: the fraction-witness
proof is the same as for
inverting $p$ in Lemma
\ref{editorial-source-lemma-p-ring-map}, replacing
each power of $p$ there by the
same power of this original $N$.
Thus it is an explicit nil-kernel
quotient, and is an isomorphism
if $\varphi$ was injective.

For a fixed prime $p$, retain
instead $n_i=p^{e_i}m_i$ and
$M=\prod_i m_i$. Over $R[1/M]$,
the global preimages above give
both
$$
x_i^{p^{e_i}}=
\frac{\varphi(A_iU_{i,p^{e_i}}+B_iE_{i,p^{e_i}})}{m_i},
\qquad
p^{e_i}x_i=\frac{\varphi(U_{i,1})}{m_i}.
$$
When $e_i=0$ these put $x_i$
itself in the image; one may
then choose exponent $1$.
For $e_i>0$ these are the
original two $p$-power conditions.
The kernel remains $I[1/M]$
and is nil. Thus the whole
localized map satisfies Lemma
\ref{editorial-source-lemma-p-ring-map}, with the
exact root exponents and all
prime-to-$p$ denominators retained.
\end{proof}